\subsection*{一步一步建立直观}

\paragraph{1. 输入和输出可以属于不同空间}

先观察三维箭头，再启用平面投影。输出只保留横纵两项，所以需要两行记录输出，却仍需三列说明三个输入方向的贡献。一个 $2$ 行 $3$ 列矩阵不是残缺的方阵，而是从三维到二维的完整规则。



\paragraph{2. 寻找消失的自由方向}

保持横纵分量不动，只改变高度。平面输出不会移动，因此竖直轴属于核。这个观察说明为什么三维输入无法由二维投影唯一恢复；丢失的是一整份自由信息，而不是计算时偶然漏掉一个数。



\paragraph{3. 把观察与维数结论分开}

每一个平面向量都可由高度为零的三维向量投过去，所以此投影是满射；它同时不是单射。维数论可以推广这种现象，但应先说明映射、定义域和值域。屏幕上的透视缩短只是相机效果，不等同于这个数学投影。
